\section{Extended Related Work}
\label{app:related-work}
\paragraph{Recovering transformations in existing representations.}
Lenc and Vedaldi \citeyearpar{lenc2015understanding} use paired linear and affine maps to study representation equivalence, including network stitching. \citet[Section~2.3 and Appendix~C.2]{williams2021shape} use orthogonal Procrustes alignment to define a distance between representations. Our independent orthogonal comparator solves this optimization separately for each prescribed transformation in the fixed, bank-whitened coordinates. The elementwise optima determine $L_{\mathrm{orth}}$; they need not satisfy the group law and are not asserted to attain $L^\star$.

\citet[arXiv v1, Section~III-C]{sun2026transformationlaws} decomposes pooled linear prediction error across realised gate-transition cells into the error of separate cell-wise linear least-squares fits and a moment-weighted cost of sharing one map. These cells distinguish the encoder's activation patterns before and after a transformation. The within-cell term can reflect both information loss and limitations of the linear family. Our comparison instead fixes the observation bank and metric, computes the cost of restricting independent linear fits to orthogonal maps, and bounds the further cost of imposing an exact group action.

\paragraph{Invariant approximation and approximate representations.}
\citet{barbieri2020approximation} construct best group-invariant approximations using orbit Gram matrices and spectral decomposition, with a budget on orbit generators in a shift-invariant setting. \citet{moore2015approximate} establish dimension-dependent obstructions to approximate representations and study compressed representations and their polar decompositions. \citet{gowers2017inverse} use averaging onto intertwiners and irreducible blocks to obtain stability results that can change the representation dimension. These tools inform our invariant-subspace construction, which retains complete real irreducible blocks under a fixed real-dimension budget. The finite-context result uses matrix concentration \citep{tropp2012user} and requires sampling assumptions beyond the fixed LLM panel.

\paragraph{Symmetry constraints, approximation and estimation.}
The regression analysis of \citet[Section~4.3, Theorem~4.9 and Corollary~4.10]{wang2023equivariance} fixes the input and output group actions and bounds approximation error by an equivariant predictor. For orthogonal output actions, the bound uses the weighted orbit variance of inversely transformed targets. Our risk formulation instead treats the output action as an optimization variable within a fixed-dimensional orthogonal representation class, and derives bounds on its optimal risk. \citet[Theorem~13]{elesedy2021generalisation} compare the expected test risks of a finite-sample minimum-norm least-squares estimate and its equivariant average when the true linear target is equivariant, under specified design and noise assumptions. Their estimation question differs from the constraint costs we compute or bound on a fixed set of observations.

\paragraph{Representational geometry and irreducibility.}
\citet{park2023linear} formalize linear concept representations and propose a metric based on causal interventions. \citet{park2025categorical} study categorical and hierarchical concept structure, and \citet{karkada2026emergent} connect language statistics to representational geometry. The statistical definition of feature irreducibility in \citet[Section~3.1]{engels2025notall} concerns whether a feature distribution can be decomposed into simpler parts. In group representation theory, irreducibility concerns the absence of proper nonzero invariant subspaces.

\paragraph{Group-operation mechanisms and competing interpretations.}
\citet{nanda2023progress} reverse engineer Fourier computations for modular addition. \citet{chughtai2023toy} recover representation-theoretic structure in networks trained on group products, and \citet{wu2025unified} develop compact performance proofs and find similar directional intervention effects in a setting where a coset-concentration explanation cannot hold. Related analyses study coset-based mechanisms \citep{stander2024cosets} and dihedral multiplication \citep{wei2026dihedral}. These studies examine computations in trained group-operation networks. Our activation edits insert the result of a composition computed externally.

\paragraph{Learning actions and measuring compositionality.}
Equivariance by Contrast \citep{schmidt2025equivariance} jointly learns latent coordinates and action representations from unlabeled action-paired observations. Its identifiability theorem assumes injective mixing, diversity of the data, exact recovery of the latent actions from their paired samples and matching conditional distributions. It recovers coordinates up to one invertible linear map, with correspondingly conjugated actions, for the points and actions in the data. Here the encoder, readout and labeled input action are fixed; we assess predictive fits in those coordinates. This observation contract differs from joint latent recovery and does not extend its theorem or transfer its guarantees to our measurements.

Homomorphism autoencoders jointly learn encoders, decoders and action representations using transition-prediction and reconstruction losses \citep{keurti2023homomorphism}. In a NeurReps workshop paper, \citet{an2026homomorphism} use representational homomorphism error as a diagnostic of compositional generalization in small Transformers on controlled SCAN-style tasks. Their fitted operators compose expression representations. Our operators predict the effects of specified group elements on contextual state representations. Equivariant architectures \citep{cohen2016group,kondor2018generalization,weiler2019general}, group crosscoders, and equivariant sparse autoencoders \citep{gorton2024group,erdogan2025group} provide symmetry-aware tools with different architectural or dictionary-learning commitments.

\paragraph{Spectral learning with a dimension constraint.}
Neural Fourier Transform \citep[Theorems~4.2 and~C.13]{koyama2024fourier} relates finite-dimensional equivariant reconstruction to the energy of irreducible components. Its formal results restrict the encoder class, including preservation of multiplicities within retained types. The nonlinear result further assumes a bi-Lipschitz invariant-kernel embedding. \citet[Sections~3--5 and Appendix~J.2]{ordonez2026inference} assume known symmetries of the marginal and conditional distributions, learn representations for Hilbert--Schmidt approximation of the conditional expectation operator, and constrain spectral truncation by irreducible dimensions and multiplicities. Their Theorem~5.1 combines operator approximation with finite-sample inference guarantees. Both settings optimize a learned representation. Our invariant-energy budget instead uses fixed, second-moment-whitened observations and a fixed real-dimension constraint. It yields bounds on the optimal exact-action risk rather than an exact minimization of that risk.

\paragraph{Probes, equivalence, and causal evidence.}
Probe selectivity requires controls for the tasks a fitted map can learn \citep{hewitt2019designing}. \citet{yao2026groupaction} emphasize the equivalence relation underlying representation claims. Their group acts on coordinate systems, while the task group here acts on input states. Exact laws survive conjugacy, whereas the reported approximate Euclidean defects depend on the selected metric. \citet{sun2026tracking} use Congruity to compare families of implementations associated with causal abstractions. Our same-element intertwiner measures compatibility of a pair of fitted operator families. Patching outcomes can also depend on corruption choices and scoring metrics \citep{zhang2024patching}. In a number-comparison task, \citet{hindupur2026comparison} find evidence for a computation using linear structure alongside curved number representations. Both observations motivate testing causal abstractions with explicit counterfactuals \citep{geiger2022inducing,geiger2025causal}.

\paragraph{Geometry, steering, and fields.}
Word relations, spatial probes, and representational convergence provide broader context \citep{mikolov2013linguistic,bolukbasi2016man,gurnee2023language,huh2024platonic}. \citet{hu2026language} study shared conceptual geometry and contextual displacements across language models. In their displacement predictor, the observed source-model displacement of a test word helps determine weights used to combine target-model training displacements. This additional information condition differs from predicting a transformation from source coordinates and an action label alone.

Activation and representation engineering \citep{turner2023activation,zou2023representation} motivate testing how representation changes affect outputs. Our constant-displacement fit uses one source-independent vector per action. Source-conditioned additive controls instead use training displacements indexed by the source state and have different information and capacity conditions. A linear operator also has an additive form, $zR_g=z+z(R_g-\Id)$, with a source-dependent displacement. These comparisons therefore concern the specified predictor classes and their information conditions; the additive form alone does not establish or exclude compositional structure.
